\documentstyle[%


righttag%


,11pt%


,amssymb%


,russian%               remove in Eng.


,izvru%                 Set izven% in Eng.


% ,izvru-ks             for brief communications


]{amsart}





%       Math definitions





\newcommand{\ind}{\operatorname{ind}}


\newcommand{\tts}[1]{}





%\hoffset=-15mm%                  For Eng.


\setcounter{page}{60}


\god{2000}


\nomer{4~(455)} %                        Remove in Eng.


%\nomer{Vol.\,44, No.\,4}%               Set in Eng.


%\str{\pageref{begin}--\pageref{end}}%  Set in Eng.


\udk{517.956}





\author{\it И.Н.\,Гурова}


% NB:   ^^ remove in Eng.


\title{О методах полудискретизации для квазилинейных уравнений


с некомпактной полугруппой}


\thanks{Работа выполнена при поддержке Российского фонда


фундаментальных исследований, грант 99-01-00333.}





\date{}


%\pagestyle{myheadings}%         Set in Eng.


\pagestyle{plain} %              Remove in Eng.


\begin{document}


%\thispagestyle{plain}%          Set in Eng.


\maketitle


\label{begin}





Метод полудискретизации для абстрактных квазилинейных уравнений в банаховом


пространстве $E$ можно описать следующим образом. Пусть в уравнении


\begin{equation}


\label{1}


x'=Ax+f(t,x)


\end{equation}


$A$ --- производящий оператор сильно непрерывной полугруппы 


$\exp\{ At \}$ линейных операторов, действующих в пространстве


$E$, а нелинейный непрерывный оператор


$f$ действует из $R^{1}\times E$ в $E$. В качестве аппроксимирующей 


полудискретизационной схемы рассмотрим уравнения


\begin{equation}


\label{2}


x'_{h} =A_{h}x_{h}+f_{h}(t,x_{h}),


\end{equation}


где $h$ --- параметр полудискретизации, операторы $A_h$ --- производящие


операторы


сильно непрерывных полугрупп $\exp\{ A_{h}t \}$ линейных операторов,


действующих в банаховых пространствах $E_{h}$, а непрерывные операторы 


$f_{h}$ действуют из $R^{1}\times E_{h}$ в $E_{h}$. Будем предполагать,


что $h \in H=\{h_n : h_n>0, \ h_n \downarrow 0 \}\cup \{ 0 \}$, при этом 


операторы $A_{0}$, $f_{0}$ отождествляются с операторами $A$, $f$, 


а пространство $E_0$ --- с пространством $E$ соответственно. Таким образом,


при $h=0$ уравнение (\ref{2}) переходит в уравнение (\ref{1}). Предположим


также, что существуют линейные ограниченные операторы вложения 


$Q_h : E_h \to E$, $h\in H\setminus \{ 0 \}$, $Q_0=I$. 


Для уравнения (\ref{1}) будем изучать


возможность приближенного нахождения решений задачи Коши или задачи о 


периодических по времени решений при помощи схемы (\ref{2}). При этом


будем говорить, что схема (\ref{2}) аппроксимирует уравнение (\ref{1}) 


в некотором функциональном пространстве ${\bold L}$, если из ограниченности


в этом пространстве последовательности 


$\{ Q_hx_h, \ h\in H\setminus \{ 0 \} \}$, где в 


качестве $x_h$ взяты решения одной из упомянутых выше задач для уравнений


(\ref{2}), следует относительная компактность этой последовательности 


в пространстве ${\bold L}$, и ее предельными точками могут быть только


решения соответствующей задачи для уравнения (\ref{1}). В случае задачи


Коши в качестве пространства ${\bold L}$ будет использоваться пространство


$C([0,d],E)$ непрерывных на некотором отрезке $[0,d]$ функций, значения


которых лежат в $E$. В задаче о периодических решениях полагаем


${\bold L}=C_T(E)$,


где $C_T(E)$ --- пространство непрерывных $T$-периодических функций со 


значениями в $E$. Под решениями соответствующих задач, следуя 


(\cite{1}, гл.\,6, \S\,23, c.\,469), понимаются решения эквивалентных


операторных уравнений, приведенных там же. 





Целью данной работы является нахождение условий, при которых схема 


(\ref{2}) 


эквивалентна такому функциональному уравнению 


\begin{equation}


\label{4}


u=F(h,u) 


\end{equation}


(в некотором пространстве ${\bold L}$),


для исследования которого может быть применена


теория вращения векторных полей в бесконечномерном пространстве 


(\cite{2}, гл.\,4, \S\,32, c.\,245).


Таким образом, задача об аппроксимации уравнения (\ref{1}) схемой (\ref{2})


сводится к изучению общими топологическими методами зависимости от параметра


решений уравнения (\ref{4}). Ранее такой подход использовался для исследования 


эллиптических \cite{3} и параболических \cite{4} уравнений.


При этом оператор $F$ оказывался вполне непрерывным. В данной работе


полугруппа $\exp\{ At \}$ предполагается лишь сильно непрерывной (это 


позволяет изучать схемы полудискретизации для квазилинейных гиперболических


уравнений). Поэтому соответствующие интегральные операторы $F$ не будут 


обладать свойством полной непрерывности.


Приведем условия, при которых соответствующий оператор будет 


уплотняющим (\cite{5}, гл.\,1, \S\,1.5, c.\,28), и поэтому для исследования 


уравнения (\ref{4}) может быть применена теория вращения уплотняющих


векторных полей (теория топологического индекса множества неподвижных


точек) (\cite{5}, гл.\,3, \S\,3.5, c.\,134).





Для построения оператора $F$ потребуются также линейные операторы


проектирования $P_h : E \to E_h$, которые удовлетворяют следующим условиям: 


\begin{equation}\label{5}


P_h Q_h =I_h,


\end{equation}


где $I_h$ --- тождественный оператор в пространстве $E_h$;


\begin{equation}\label{6}


Q_h P_h x \to x \ \; \text{для  всех}\ \; x\in E.


\end{equation}


Всюду ниже считаем, что $\| P_h \|$ и $\| Q_h \|$ равномерно 


ограничены.





Перейдем теперь к условиям на аппроксимирующие операторы $A_h$ и $f_h$.


Будем предполагать, что полугруппы $\exp\{A_h t\}$ аппроксимируют полугруппу


$\exp\{At\}$, т.\,е. выполняется условие


{\it





$\mathrm A_1)$ для любого $x\in E$


$$


Q_h \exp\{A_h t\} P_h x \longrightarrow \exp\{At\}x


\ \; \text{при}\ \; h\to 0


$$ 


равномерно по $t$ из любого ограниченного отрезка.


}





Для того чтобы сформулировать условия на операторы $f_h$, рассмотрим 


оператор $\varphi : H\times \bold R^1\times E \to E$, задаваемый формулой 


$$


\varphi(h,t,x)= Q_h f_h (t,P_h x).


$$





Напомним, что мерой некомпактности Хаусдорфа 


(\cite{5}, c.\,7) ограниченного


множества $\Omega\in E$ называют число $\chi (\Omega )$,


определяемое равенством


$$


\chi (\Omega)= \inf \{\varepsilon :


\Omega\ \; \text{имеет конечную $\varepsilon$-сеть}\}.


$$ 





Ниже предполагаем выполненными следующие условия:


\begin{itemize}


\item[$\mathrm A_2)$] {\it оператор $\varphi$ непрерывен по совокупности 


переменных и ограничен на ограниченных множествах}; 


\item[$\mathrm A_3)$] {\it существует такая константа $k$, что 


$$


\chi\Big(\bigcup\limits_{h\in H}\, \varphi(h,t,\Omega)\Big)


\le k\chi(\Omega)


$$


для любого ограниченного множества $\Omega \in E$}.


\end{itemize}





Для уравнения (\ref{1}) рассмотрим задачу Коши с начальным условием 


\begin{equation}\label{7}


x(0) =x^0.


\end{equation}





В качестве аппроксимации (\ref{7}) возьмем для уравнения (\ref{2}) 


начальное условие


\begin{equation}\label{8}


x_h (0) =P_h x^0.


\end{equation}





Оператор $F: H\times C([0,d],E)\to C([0,d],E)$ определим с помощью формулы


\begin{equation}\label{9}


F(h,u)(t)=Q_h \exp\{A_h t\}P_h x^0


+\int_0^t Q_h \exp\{A_h (t-s)\}P_h \varphi(h,s,P_h u(s))ds.


\end{equation}





Далее используется мера некомпактности $\psi$ \cite{6}, определяемая 


на ограниченных подмножествах $\Omega$ пространства $C([0,d],E)$


формулой


$$


\psi (\Omega)= \max_{D\subseteq\Omega}(\sup_{t} e^{-Lt} \chi (D(t)),


\quad \lim_{\delta \to 0}


\sup_{u\in D}\max_{0\le\tau\le\delta }\| u(t)-u(t+\tau)\|_E),


$$


при этом максимум в ee правой части берется по всем счетным подмножествам


$D$ множества $\Omega$, а неотрицательная константа $L$ носит


технический характер и выбирается в ходе доказательства.


Значения меры некомпактности $\psi$ лежат в конусе 


$\bold R^2_+$. В смысле порядка, индуцированного этим конусом,


и понимаются ниже


неравенства, связанные с мерой некомпактности.





Как показано в \cite{6}, мера некомпактности $\psi$ обладает следующими


свойствами:


\begin{itemize}


\item[1)] монотонность, т.\,е. из включения $\Omega_1\subseteq \Omega_2$ 


следует неравенство $\psi(\Omega_1)\leq \psi(\Omega_2);$ 


\item[2)] инвариантность относительно присоединения компактного


множества, т.\,е.


для любого компактного $K\subset C([0,d],E)$ и ограниченного 


$\Omega\subset C([0,d],E)$ выполнено равенство


$$


\psi (\Omega\cup K)=\psi (\Omega);


$$


\item[3)] правильность, т.\,е. равенство $\psi (\Omega)=0$ эквивалентно 


относительной


компактности множества $\Omega$.


\end{itemize}





Поэтому (\cite{5}, c.\,134) для операторов, 


уплотняющих относительно меры 


некомпактности $\psi$, справедлива теория топологического индекса множества


неподвижных точек, построенная там же. 


Напомним, что оператор 


$G$ называется уплотняющим относительно меры некомпактности $\psi$, 


если он непрерывен, и из неравенства


$\psi (G\Omega)\geq \psi (\Omega)$


вытекает относительная компактность множества $\Omega$. 


Говорят, что $F$ уплотняет по совокупности переменных относительно меры


некомпактности $\psi$, если он непрерывен, и из неравенства


$\psi (F(H\times \Omega))\geq \psi (\Omega)$ вытекает


относительная компактность множества $\Omega$. 


 


\begin{thm}\label{th1}


Пусть выполнены условия $(\ref{5})$, $(\ref{6})$ и предположения


$\mathrm A_1)$--$\mathrm A_3)$.


Тогда оператор~$F$, задаваемый формулой $(\ref{9})$,


уплотняет по совокупности переменных относительно


меры некомпактности $\psi$. Соотношение между неподвижными точками


$u^h$ оператора $F(h,\cdot)$ и решениями 


$x_h$ задачи $(\ref{2})$, $(\ref{8})$ задается равенствами


\begin{equation}\label{9'}


P_h u^h = x_h, \qquad Q_h x_h =u^h.


\end{equation}


\end{thm}





Из полученного результата следует, что для задачи Коши


схема (\ref{2}) аппроксимирует уравнение (\ref{1}).





Перейдем теперь к условиям, обеспечивающим непустоту множества решений


уравнений (\ref{2}).


Обозначим через $\Sigma_h [0,\tau ]$ 


множество решений задачи (\ref{2}), (\ref{8}),


определенных на отрезке $[0,\tau ]$. Ниже будем предполагать, что 


для любого $\tau < d$ все решения задачи (\ref{1}), (\ref{7}), определенные


на отрезке $[0,\tau ],$ могут быть продолжены на отрезок $[0,d]$. Это 


предположение запишем в виде





$\mathrm A_4)$\; $\Sigma_0 [0,d] \big|_{[0,\tau ]} = \Sigma_0 [0, \tau].$ 





Как показано в \cite{6}, если выполнено предположение $\mathrm A_4)$


и множество $\Sigma_0 [0,d]$ ограничено, то 


$$


\ind (\Sigma_0 [0,d], F(0,\cdot)) = 1.


$$ 





Применяя теперь теорему \ref{th1} и аналог теоремы 54.1 (\cite{2}, 


c.\,458) для уплотняющих операторов, получим следующий результат.





\begin{thm}\label{th2} 


Пусть выполнены условия $(\ref{5})$, $(\ref{6})$, предположения


$\mathrm A_1)$--$\mathrm A_4)$ и множество $\Sigma_0 [0,d]$ ограничено. 


Тогда при достаточно малых $h$ все решения задач 


$(\ref{2})$, $(\ref{8})$ определены на отрезке $[0,d]$ и отображение


$h\mapsto Q_h \Sigma_h [0,d]$ полунепрерывно сверху. 


\end{thm} 





Из последней теоремы, в частности, следует: если задача 


(\ref{1}), (\ref{7}) имеет на отрезке $[0,d]$ единственное решение $x^0$,


то при достаточно малых $h$ все решения задач (\ref{2}), (\ref{8})


определены на отрезке $[0,d]$ и


\begin{equation}\label{9''}


Q_hx_h \to x^0 \ \; \text{при} \ \; h\to 0,


\end{equation}


где $x_h\in \Sigma_h [0,d]$.





Перейдем теперь к задаче о периодических решениях уравнения (\ref{1}).


При изучении этой задачи 


$f$ и $f_h$ предполагаются $T$-периодичными


по первой переменной. Кроме того, считаем


пространство $E$ сепарабельным. 





Определим оператор $F : H\times C_T (E) \to C_T (E)$


формулой


\begin{multline}\label{10}


F(h,u)(t)=Q_h \exp\{A_ht\} (I- \exp\{A_hT\})^{-1} 


\int_0^T \exp\{A_h(T-s)\}P_h\varphi(h,s,P_hu(s))ds +\\


%


+\int_0^t Q_h\exp\{A_h(t-s)\}P_h\varphi(h,s,P_hu(s))ds.


\end{multline}


Отметим, что условие $1\notin \sigma(\exp\{A_ht \})$, необходимое для 


корректности последней формулы, носит формальный характер. Его всегда можно 


предполагать выполненным. В противном случае нужно в правой части уравнения


(\ref{2}) вычесть и прибавить член $Lx_h$, где $L$ --- достаточно большая


положительная константа, и затем рассмотреть операторы $A_hx_h -Lx_h$ и 


$f_h(t, x_h)+Lx_h$ в качестве новых операторов $A_h$ и $f_h$. Нетрудно 


показать, что если для первоначальных операторов $A_h$ и $f_h$ были


выполнены предположения $\mathrm A_1)$--$\mathrm A_3)$ и накладываемое


ниже требование


$\mathrm A_5)$, то предположения $\mathrm A_1)$--$\mathrm A_5)$


выполнены и для переопределенных операторов. 





В пространстве $C_T (E)$ будем рассматривать меру некомпактности 


$\nu$ \cite{7}, задаваемую формулой


$$


\nu (\Omega) =\big(\chi (\Omega)(\cdot), \ \lim_{\delta \to 0}


\sup_{u\in D}\max_{0\le\tau\le\delta }\| u(t)-u(t+\tau)\|_E\big).


$$


Значения меры некомпактности $\nu$ лежат в конусе


$\bold K = K_T\times \bold R_+^1$, 


где $K_T$ --- конус измеримых, ограниченных почти всюду $T$-периодических


функций. Мера некомпактности $\nu$ так же, как и мера некомпактности $\psi$,


обладает свойствами 1)--3). Поэтому для уплотняющих относительно $\nu$ 


операторов тоже справедлива теория топологического индекса неподвижных


точек. 





Кроме $\mathrm A_1)$--$\mathrm A_3)$,


при изучении задачи о периодических 


решениях потребуется еще предположение


 


$\mathrm A_5)$ {\it существует константа $\gamma > k$ такая, что}


$$


\chi\Big(\bigcup\limits_{h\in H}\, \ Q_h \exp\{A_ht \} P_h B_E (0,1)


\Big)\le e^{-\gamma t}.


$$





Наконец, предположим, что для операторов $Q_h$, $P_h$ выполнено неравенство


\begin{equation}\label{p''}


\chi \Big(\bigcup_{h\in H}Q_hP_hB(0,1)\Big)\le 1.


\end{equation}





\begin{thm}\label{th3}


Пусть выполнены условия $(\ref{5})$, $(\ref{6})$, $(\ref{p''})$


и предположения


$\mathrm A_1)$--$\mathrm A_3)$, $\mathrm A_5)$.


Тогда оператор $F$, задаваемый формулой $(\ref{10})$,


уплотняет по совокупности переменных относительно


меры некомпактности $\nu$. Соотношение между неподвижными точками


$u^h$ оператора $F(h,\cdot)$ и $T$-периодическими решениями 


$x_h$ уравнения $(\ref{2})$ задается равенствами $(\ref{9'})$.


\end{thm}





Применяя, как и выше, аналог теоремы 54.1 (\cite{2}, c.\,458), 


получим следующий результат.





\begin{thm} Пусть выполнены условия $(\ref{5})$, $(\ref{6})$ и предположения


$\mathrm A_1)$--$\mathrm A_4)$. Предположим, что $x^0$ ---


$T$-периодическое решение


уравнения $(\ref{1})$, для которого 


\begin{equation}\label{11}


\ind (x^0, F(0,\cdot))\ne 0.


\end{equation} 


Тогда для достаточно малых $h$ уравнения $(\ref{2})$ имеют $T$-периодические


решения $x_h$, для которых выполнено соотношение $(\ref{9''})$.


\end{thm} 





Неравенство (\ref{11}) выполнено, например, если оператор $F(0,\cdot)$ 


дифференцируем в точке $x^0$ и 


\begin{equation}\label{12}


1\notin \sigma(F'_x(0, x^0)).


\end{equation} 


Отметим, что в случае задачи Коши последнее соотношение всегда выполнено,


т.\,к. в силу теоремы


единственности линеаризованное уравнение с нулевым начальным условием не 


может иметь ненулевых решений. 


 


Оценим теперь скорость сходимости $Q_h x_h$ к $x^0$ или 


в силу соотношений (\ref{9'}) $u^h$ к $x^0$. При этом будем предполагать,


что операторы $F(h, \cdot)$ дифференцируемы по второй переменной в точке


$x^0$ равномерно относительно первой, т.\,е.


\begin{equation}\label{13}


F(h,x_0+y)-F(h,x_0)=F'_x(h,x^0)y+\omega(h,y),


\end{equation}


где $\|\omega(h,y)\|/\|y\|\to 0$ при $\|y\|\to 0$ равномерно относительно


$h$. Для выполнения последнего достаточно потребовать, чтобы операторы


$\psi(h,t,x)$ были дифференцируемы по третьей переменной в точках $x_0(t)$


равномерно относительно $h$, $t$. Кроме этого, будем предполагать, что


выполнено соотношение (\ref{12}). Так как оператор $F'_x(h,x^0)y$ уплотняет 


по совокупности переменных $h$, $y$, то в силу 


леммы 2.7.7 (\cite{5}, c.\,98) 


$1\notin \sigma(F'_x(h, x^0))$ при достаточно малых $h$ и 


$\|(I-F'_x(h, x^0))^{-1}\|$ равномерно ограничены. Поэтому из соотношения


(\ref{13}) вытекает эквивалентность $\|u^h - x^0\|$ и 


$\|F(h, x^0)-F(0, x^0)\|$. Таким образом, сходимость $Q_hx_h$ к $x^0$ 


эквивалентна сходимости аппроксимаций на отыскиваемом решении.





В заключение остановимся на некоторых примерах выполнения условий 


$\mathrm A_1)$--$\mathrm A_5)$.





Пусть $A_h=\wh{A}_h+C_h$, причем 


операторы $\wh{A}_h$ порождают в пространствах $E_h$ полугруппы


$\exp\{\wh{A}_ht\},$ удовлетворяющие при некотором $\gamma>0$ оценке 


\begin{equation}\label{14}


\|\exp\{\wh{A}_ht\}\|_{E_h}\leq e^{-\gamma t},


\end{equation} 


а операторы $C_h$ таковы, что $Q_hC_hP_hx$ вполне непрерывен по совокупности


переменных $h\in H$, $x\in E$. Пусть также $\|Q_h\|,\, \|P_h\|\leq 1$ и


существует множество $D$, плотное в $E$,


содержащееся при некотором $\lambda\in \bold C$ в


$(\lambda-A)^{-1}E$ и такое, что 


$$


\|P_hAx-A_hP_hx\|_{E_h}\to 0\ \; \text{при}\ \; h\to 0, \ \; x\in D.


$$ 


Тогда для операторов $A_h$ выполняются требования


$\mathrm A_1)$ и $\mathrm A_5)$. Для 


выполнения оценки (\ref{14}) достаточно в силу теоремы Хилле--Иосиды


(\cite{8}, c.\,343), чтобы для 


вещественных $\lambda \geq -\gamma$ выполнялось неравенство


$$


\|(\lambda-\wh{A}_h)^{-1}\|_{E_h}\leq \frac {1}{\lambda+\gamma}.


$$


 


Отметим, что перечисленные условия реализуются для аппроксимаций оператора, 


возникающего при описании передаточных линий, описываемых квазилинейными


телеграфными уравнениями \cite{9}, если в первом уравнении производная 


по пространственной переменной аппроксимируется разностью вперед, а


во втором --- разностью назад. В качестве $P_h$ нужно взять


проекторы Стеклова


по отрезкам, образуемым сеткой, а в качестве $Q_h$ --- оператор,


естественным образом сопоставляющий сеточной функции ступенчатую. 





Если нелинейность $f$ в уравнении (\ref{1}) непрерывна, ограничена


на ограниченных множествах, удовлетворяет


условию


$$


\chi (f(t,\Omega))\leq k\chi (\Omega)


$$


и, кроме того, $\|Q_hP_h\|\leq 1$,


то оператор $\varphi$, порожденный оператором


$$f_h(t,x_h)=P_hf(t,Q_hx_h),


$$


удовлетворяет требованиям


$\mathrm A_2)$, $\mathrm A_3)$. Если дополнительно выполнено


неравенство


$$


\|f(t,x)\|\le a(1+\|x\|),


$$ то требование $\mathrm A_4)$ выполнено


для любого $d>0$. 





\begin{thebibliography}{9}





\bibitem{1}


Красносельский М.А., Забрейко П.П., Пустыльник Е.И., Соболевский П.Е. 


{\em Интегральные операторы в пространствах суммируемых функций}.


-- М.: Наука, 1966. -- 499 с.


\bibitem{2}


Красносельский М.А., Забрейко П.П. 


{\em Геометрические методы нелинейного анализа}.


-- М.: Наука, 1975. -- 510 с.


\bibitem{3}


Гурова И.Н.


{\em Об одном топологическом методе исследования разностных 


схем }//


ДАН СССР. -- 1979. -- Т.\,248. -- \No\,1. -- С.\,25--28.


\bibitem{4}


Гурова И.Н., Каменский М.И.


{\em О методе полудискретизации в задаче о периодических решениях


квазилинейных автономных параболических уравнений} //


Дифференц. уравнения. -- 1996. -- Т.\,32. -- \No\,1. -- С.\,101--106.


\bibitem{5}


Ахмеров Р.Р., Каменский М.И., Родкина А.Е., Потапов А.С., Садовский Б.Н.


{\em Меры некомпактности и уплотняющие операторы}.


-- Новосибирск: Наука, 1986. -- 282 c.


\bibitem{6}


Couchouron J.-F., Kamenskii M.I.


{\em A unified topological point of view for integro-differential 


inclusions} //


Different. Inclusions and Optim. Control.


Lect. Notes in Nonlinear Anal. -- 1998. -- V.\,2. -- P.\,123--137.


\bibitem{7}


Kamenskii M.I.


{\em Mesures de non compacit\'e dans l'espace des fonctions


continues et leurs applications aux \'equations et aux inclusions 


diff\'erentielles} //


S\'emin. Initiation \`a l'Analyse, 


Publ. Math. de l'Universit\'e Pierre et Marie Curie.


-- 1996/97. -- V.\,120. -- P.\,11--13. 


\bibitem{8}


Иосида К. {\em Функциональный анализ}. -- М.: Мир, 1967. -- 624 c. 


\bibitem{9}


Ceron S.S., Lopes O. {\em $\alpha$--contractions and attractors for


dissipative semilinear hyperbolic equations and systems} //


Ann. Math. Pura Appl.


-- 1991. -- V.\,160. -- \No\,4. -- P.\,193--206.





\end{thebibliography}





\vskip 2cm





\post{\it Воронежский государственный}{\it Поступила}


\post{\it технический университет}{30.04.1999}





\label{end}


\end{document}


